\documentclass[12pt,a4paper]{article}

% --- Basic Packages (Strict ASCII, no unicode hooks) ---
\usepackage[utf8]{inputenc} % Keep source code clean, but we stick to pure LaTeX commands for math/symbols
\usepackage[T1]{fontenc}
\usepackage{amsmath}
\usepackage{amsfonts}
\usepackage{amssymb}
\usepackage{geometry}
\usepackage{listings}

\geometry{margin=1in}

\title{\textbf{Project KHAL-VM: Architecture Decision Records}}
\author{Kirill \and HAL-9000}
\date{July 2026}

\begin{document}

\maketitle

\thispagestyle{empty}
\newpage
\pagestyle{plain}

\section{Introduction and Development Strategy}
Project KHAL-VM is a research functional virtual machine designed iteratively from scratch. The development roadmap follows three distinct evolutionary steps:
\begin{enumerate}
    \item \textbf{Sandbox Phase:} Rapid prototyping of execution core, Instruction Set Architecture (ISA), and frontend syntax within the Emacs environment using light scripting engines (F\# or Chez Scheme).
    \item \textbf{Metropolis Phase:} Complete transition of the proven architecture to pure C++ based on Bjarne Stroustrup's principles (RAII resource management) and Alexander Stepanov's paradigms (mathematical rigor and generalized algorithms).
\end{enumerate}

\section{Core Architectural Decisions (ADR)}

\subsection{ADR-001: Register-Based Execution Model}
\textbf{Context:} Traditional managed virtual machines (e.g., JVM, CLR) utilize stack-based bytecode for compiler simplicity. This introduces substantial memory thrashing (via continuous PUSH/POP cycles) and high instruction density. \\
\textbf{Decision:} KHAL-VM implements a pure \textbf{register-based bytecode format}, drawing inspiration from LuaJIT and Erlang BEAM. \\
\textbf{Consequences:} The instruction count per execution block drops by 2--3x. Dispatch loop pressure on the interpreter is minimized. The trade-off is shifted toward a more complex register allocation mechanism.

\subsection{ADR-002: Hybrid Memory Topography}
\textbf{Context:} Physical hardware operates on combined stack/register logic. High-frequency functional invocations require lightweight state transition. \\
\textbf{Decision:} KHAL-VM bifurcates its execution storage:
\begin{enumerate}
    \item \textbf{Register File (256 slots):} Utilized for localized evaluation and pure, static functions. Arguments are passed with zero-cost overhead via register window overlapping.
    \item \textbf{Call Stack (1024 frames):} Dedicated strictly to control flow tracking, Instruction Pointer (IP) state preservation, and execution context storage.
\end{enumerate}

\subsection{ADR-003: Tagged Hardware Architecture}
\textbf{Context:} Offloading type checking completely to static compilation stages complicates the frontend and limits runtime flexibility. \\
\textbf{Decision:} Mirroring vintage Lisp-machine designs, every KHAL-VM register contains two fixed components: \texttt{Tag} (1 byte, type descriptor) and \texttt{Bits} (8 bytes, raw payload data). \\
\textbf{Consequences:} Attempting operations on incompatible registers triggers an immediate hardware-level interrupt. The frontend compiler is completely stripped of heavy type-checking logic.

\subsection{ADR-004: 64-bit Mathematical Payload Layout}
\textbf{Context:} Efficient bitwise arithmetic is required for basic numerical types. \\
\textbf{Decision:} Three core tags are established for the initial mathematical unit:
\begin{itemize}
    \item \texttt{TAG\_INT} (0x01): Standard 64-bit signed integer storage.
    \item \texttt{TAG\_FIXED} (0x02): Fixed-point layout (32-bit integer, 32-bit fractional partition) utilizing integer ALU speeds.
    \item \texttt{TAG\_RATIONAL} (0x03): Strict fractional precision. The payload is divided precisely in half: upper 32 bits store the \textbf{Numerator}, lower 32 bits store the \textbf{Denominator}. 
\end{itemize}
Extraction is executed on the host system via basic bitwise masking and shifting operations:
\begin{align}
    \text{Denominator} &= \text{Bits} \ \& \ \text{0xFFFFFFFF} \\
    \text{Numerator} &= \text{Bits} \ >> \ 32
\end{align}

\subsection{ADR-005: Wirthian Lisp Frontend Compiler}
\textbf{Context:} A parser is required to map standard S-expressions directly into register instructions. \\
\textbf{Decision:} KHAL-VM integrates a prefix Lisp syntax but compiles it following Niklaus Wirth's constraints --- strictly in \textbf{one single pass} without constructing an Abstract Syntax Tree (AST) in heap memory. \\
\textbf{Consequences:} For example, the expression \texttt{(* (+ a b) c)} is parsed linearly from left to right. Upon reaching a closing parenthesis, the compiler immediately emits specific register operations:
\begin{lstlisting}[basicstyle=\ttfamily]
1. ADD R1, a, b
2. MUL R2, R1, c
\end{lstlisting}
Constraint: Strict declaration-before-use policy.



\section{Sandbox Execution: First Operational Milestone}

\subsection{ADR-007: Multi-Engine Directory Isolation and Scripting Strategy}
\textbf{Context:} Rapid architectural validation requires evaluating syntax and execution mechanics simultaneously across disparate runtime paradigms without early infrastructure overhead. \\
\textbf{Decision:} The repository layout is partitioned into isolated, self-contained subdirectories to decouple testing environments completely:
\begin{itemize}
    \item \texttt{khal\_vm\_f/}: Dedicated to F\# lightweight scripting (\texttt{.fsx}) via the interactive compiler (\texttt{dotnet fsi}).
    \item \texttt{khal\_vm\_s/}: Dedicated to Chez Scheme scripting (\texttt{.ss}) executed linearly via the native engine (\texttt{scheme --script}).
\end{itemize}
\textbf{Consequences:} This eliminates project configuration files (\texttt{.fsproj} or scheme libraries) during the initial R\&D phase, allowing rapid iterative coding directly within Emacs buffers. Transition to rigid build systems is deferred until multi-file cohesion becomes required.

\subsection{ADR-008: Instruction 0x0A (PRINT) and String Pool Emulation}
\textbf{Context:} A register machine must validate its instruction decoding loop using a predictable, verifiable operational primitive before scaling the Instruction Set Architecture (ISA). \\
\textbf{Decision:} KHAL-VM establishes opcode \texttt{0x0A} (\texttt{OP\_PRINT}) as the initial functional baseline. String variables are decoupled from the core register file to maintain a strict, non-OOP structural topography:
\begin{enumerate}
    \item A localized \texttt{String Pool} mimics static memory read segments.
    \item Register \texttt{R0} is loaded with a pure 64-bit integer acting as an indirect address offset (index) pointing into the pool.
    \item The interpreter loop fetches the instruction, extracts the register reference, reads the index parameter, and pipe-routes the data to the host stdout.
\end{enumerate}
\textbf{Consequences:} The primary execution mechanism (Fetch-Decode-Execute cycle) is successfully instantiated in both F\# and Scheme sandboxes, printing the historical "Hello World" paradigm with zero object allocations.


\begin{lstlisting}[basicstyle=\ttfamily\small, frame=single, title=Verified F\# Sandbox Execution Log]
[LAUNCH]  KHAL-VM (F#) activated in isolated environment.
[VM OUTPUT]  Hello, World! (KHAL-VM F#)
[READY]   Execution completed. IP = 0x01.

val OP_PRINT: int = 10
val stringPool: string array = [|"Hello, World! (KHAL-VM F#)"|]
val registers: int array = [|0; 0; 0; 0; 0; ...|]
val mutable IP: int = 1
val bytecode: (int * int) array = [|(10, 0)|]
\end{lstlisting}


\begin{lstlisting}[basicstyle=\ttfamily\small, frame=single, title=Verified Chez Scheme Sandbox Execution Log]
=> #<void>

[LAUNCH]  KHAL-VM (Scheme) activated in isolated environment.
[VM OUTPUT]  Hello, World! (KHAL-VM Scheme)
[READY]   Execution completed. IP = 0x1.
\end{lstlisting}

\end{document}

